\chapter{Conclusion and Use Cases}
\label{ch:conclusion}

\section{Conclusion}

This work set out to make natural-language access to a university ERP safe enough to try on real users.
The answer NatQuery gives is architectural: a language model may propose SQL, but it is never trusted to
run it. Every query is parsed, checked against the asker's role view, rewritten with row filters and
masks, regenerated, recorded before it runs, and executed read-only, with PostgreSQL row-level security
repeating the row scopes underneath.

The evidence supports three claims. First, the security layer is almost free: on BIRD Mini-Dev it
rejected 2\% of first attempts, turned no correct answer into a wrong one, and added 4\,ms to a request
that takes a second, while an open 120-billion-parameter model reached 45.6\% execution accuracy, ahead of
GPT-4 given the same hints. Second, it holds under attack: 560 hostile-SQL runs and 30 natural-language
attacks leaked nothing, where an unprotected text-to-SQL leaked on 7 of 8 comparable attacks. Third, it
scales: on a synthetic university ERP with 87 tables, 549 columns and 10 roles, the access layer returned
exactly the right rows on every question, and the system answered 28 of 40 correctly with all five
access-control traps handled safely.

Equally important is what the work does not show. Accuracy on the ERP depends on the model and on
descriptions of the schema, and the final figure comes from improvements designed after reading the
first run's failures. The attack list was written by the authors of the defences. The real ERP's schema
and policy remain unknown. These are the first things to address before any real deployment.

\section{Use Cases}

Table~\ref{tab:usecases} lists the intended uses at a university. They are possible uses that the test
bed exercises, not deployments.

\begin{table}[htbp]
\centering
\caption{Intended use cases.}
\label{tab:usecases}
\small
\renewcommand{\arraystretch}{1.15}
\begin{tabularx}{\textwidth}{L{2.6cm} Y Y}
\toprule
\hdr \textbf{User} & \textbf{Example question} & \textbf{What the policy guarantees} \\
\midrule
Student & ``In which courses is my attendance below 95 percent?'' & Own rows only; no other student's data; secrets never visible \\
Faculty & ``Which of my advisees have an unpaid library fine?'' & Own courses and advisees; own salary; no other staff pay \\
Head of department & ``Who are the five students in my department with the highest CGPA?'' & Department scope; no other department; masked fields stay masked \\
Accounts section & ``How many fee payments failed this semester?'' & Finance tables for the role; no academic detail beyond identification \\
Library staff & ``Which books are overdue by more than a month?'' & Library tables and the minimum student fields needed \\
Hall office & ``How many residents of my hall have unpaid dues?'' & One hall only \\
Placement cell & ``For companies with over 90 applications, what percentage received offers?'' & Placement tables only \\
Administrator & ``Show recent query activity and verify the audit chain.'' & Audit commands; permanent secrets stay hidden even for admin \\
Help desk & Common ERP help-desk questions answered without a human & The same row scoping as the logged-in user \\
Other institutions & Any PostgreSQL system with roles and personal records & Policy file and schema description replace the ERP-specific parts \\
\bottomrule
\end{tabularx}
\end{table}

\section{Before a Real Deployment}

\begin{table}[htbp]
\centering
\caption{Deployment checklist.}
\label{tab:checklist}
\small
\renewcommand{\arraystretch}{1.15}
\begin{tabularx}{\textwidth}{L{4.2cm} Y}
\toprule
\hdr \textbf{Step} & \textbf{Why} \\
\midrule
Write the policy against the real schema and review every role's view & The policy is only as good as its author; \code{natquery policy check} validates it \\
Use a dedicated SELECT-only database account and row-level security & A second line of defence beneath the enforcer \\
Decide where the model runs & The question and the role-visible schema leave the machine; an on-premises model keeps them in; row data never leaves \\
Ship the audit log off the machine and set the audit key & The chain is tamper-evident, not tamper-proof \\
Add a notice when a row filter changed an answer & Silent scoping can mislead a user, although nothing leaks \\
Run an outside red team & The attack list was written by the authors \\
Add per-device throttling and shared session storage & Per-user lockout can be abused; in-memory limits suit one process \\
\bottomrule
\end{tabularx}
\end{table}

\section{Limitations}

\begin{table}[htbp]
\centering
\caption{Known limitations.}
\label{tab:limits}
\small
\renewcommand{\arraystretch}{1.15}
\begin{tabularx}{\textwidth}{L{3.8cm} Y}
\toprule
\hdr \textbf{Limitation} & \textbf{Detail} \\
\midrule
Silent scoping & A row filter changes totals without telling the user: ``how many employees work here'' returns 1 for an employee; the model cannot see filters, so some ``my'' questions are judged unanswerable \\
No minimum group size & A role allowed to aggregate can learn statistics about rows it cannot see \\
Heuristic input guard & A speed bump and audit signal; the enforcer is the boundary \\
Single process & Sessions and rate limits are in memory; the audit chain is tamper-evident only \\
One role per query & By design, for simplicity of reasoning and audit; no union of roles \\
PostgreSQL only & Other engines need a backend class and a dialect name \\
Multi-schema naming & Tables that exist in several schemas must be written in full by the model, which costs repair attempts \\
Row-level security covers rows only & Column hiding and masking live in the enforcer \\
Evaluation scope & Synthetic ERP, one model, one run each, partly tuned on the benchmark questions \\
\bottomrule
\end{tabularx}
\end{table}

\section{Future Work}

\begin{table}[htbp]
\centering
\caption{Future work.}
\label{tab:future}
\small
\renewcommand{\arraystretch}{1.15}
\begin{tabularx}{\textwidth}{L{4cm} Y}
\toprule
\hdr \textbf{Direction} & \textbf{Detail} \\
\midrule
Real ERP trial & Connect to the institute ERP (or a copy), collect real help-desk questions and write the real policy \\
Scope-aware answers & Tell the user when a filter limited the result; let the model know which tables are already scoped to the user \\
Aggregate protection & Minimum group sizes and noise for inference-prone aggregates \\
Stronger models & Re-run the ERP and BIRD tracks with a larger or on-premises model and higher reasoning effort \\
Outside security review & Independent red team and a published attack corpus \\
Production layer & Network API, single sign-on (LDAP or OIDC), shared limiter and session store, concurrency testing \\
More engines & Backends for MySQL and Oracle through the backend interface \\
Learning from feedback & Use rejected and repaired queries to improve descriptions and value hints automatically \\
\bottomrule
\end{tabularx}
\end{table}

\FloatBarrier
